% Fichier complété par tools/complete_missing_corrections.py.
% Source : SLCI/SLCI-07-Ordre12/541_Cours/541_Cours.tex
% Statut : rédigé (5 question(s)).
\begin{corrigebox}[Corrigé rédigé]
\CorrigeQuestion{1}
$H(p)=K/(1+\tau p)$.
\CorrigeQuestion{2}
Pour un échelon unitaire et $K=0,5$ :
            $s(t)=0,5(1-e^{-t/\tau})u(t)$. La tangente à l'origine coupe la valeur finale
            en $t=\tau$.
\CorrigeQuestion{3}
L'écart est $\varepsilon(t)=1-s(t)$; sa valeur finale vaut
            $\varepsilon_\infty=1-K=0,5$ pour une chaîne directe seule.
\CorrigeQuestion{4}
$t_{5\%}\simeq3\tau$.
\CorrigeQuestion{5}
Méthodes : (1) tangente à l'origine, donnant $\tau$; (2) lecture du temps
            où la réponse atteint $63\%$ de sa variation finale. Le gain est le rapport
            valeur finale/valeur de l'échelon.
\end{corrigebox}
